\documentclass[12pt,aspectratio=169,ignorenonframetext]{ltxbeamer}

\usepackage[spanish,es-nodecimaldot,es-tabla]{babel}
\usepackage{animate}
\usepackage{multicol}
\usepackage{multimedia}
\usepackage{pgfplots}
\usepackage{pgfplotstable}
\sisetup{
    uncertainty-mode=separate,
}
\pgfplotsset{
    compat=1.18,
}
\pgfplotstableset{
    col sep=comma,
}

\newcommand{\beamer}{\textsc{beamer}}

\title{Presentaciones {\lmr\mdseries\beamer}}
\subtitle{Curso de {\lmr\LaTeX}}
\author{Joar Buitrago\texorpdfstring{\thanks{jebuitragoc@unal.edu.co}}{ <jebuitragoc@unal.edu.co>}}
\date{}

\addbibresource{bibliography.bib}



\begin{document}
\begin{frame}
\titlepage
\end{frame}



\section{¿Por qué {\lmr\beamer}?}

\begin{frame}
\frametitle{El reto de comunicar ciencia}
\begin{itemize}
    \item Las herramientas tradicionales (PowerPoint/Keynote) suelen fallar en el manejo de tipografía matemática compleja.
    \item La desvinculación entre el «Artículo» y la «Presentación» genera doble trabajo.
    \item {\lmr\beamer} permite reusar el código exacto de su tesis o artículo en sus diapositivas.
\end{itemize}
\vfill
\centering
\textit{«Una presentación científica debe ser tan rigurosa como el manuscrito que sustenta.»}
\end{frame}



\begin{frame}
\frametitle{Ventajas de {\lmr\beamer}}
\begin{itemize}[<+->]
    \item Mismas fuentes y calidad que en {\lmr\LaTeX}. (O fuentes personalizadas con \texttt{fontspec} y \texttt{unicode-math})
    \item Control exacto de la aparición de datos.
    \item Navegación interactiva inmediata.
    \item Un PDF que se ve igual en cualquier computadora, sin problemas de fuentes o versiones.
    \item Animaciones programables desde el código.
\end{itemize}
\end{frame}



\section{Anatomía de un Proyecto {\lmr\beamer}}

\begin{frame}[fragile]
\frametitle{El preámbulo mínimo}
Para iniciar, debemos cambiar la clase del documento a {\lmr\beamer}.

\inputminted{latex}{resources/minimum.tex}
\end{frame}




\begin{frame}
\centering
\includegraphics[width=0.7\linewidth]{resources/build/minimum.pdf}
\end{frame}




\begin{frame}[fragile]
\frametitle{El entorno \texttt{frame}}
Cada diapositiva es su propio entorno (\texttt{frame}).

\begin{description}
    \item[\mintinline{latex}{\frametitle}] Define el encabezado.
    \item[\mintinline{latex}{\framesubtitle}] Información secundaria.
\end{description}
\end{frame}





\begin{frame}[fragile]
\frametitle{El entorno \texttt{frame}}

\only<1>{\inputminted[firstline=4,lastline=7]{latex}{resources/frame.tex}}
\only<2>{
    \centering
    \includegraphics[width=0.5\linewidth]{resources/build/frame.pdf}\par
}
\end{frame}





\section{Organización de Contenidos}

\begin{frame}[fragile]
\frametitle{Uso efectivo de Bloques}
Los bloques agrupan ideas lógicas. {\lmr\beamer} ofrece tres tipos estándar por defecto:

\only<1>{\inputminted[firstline=5,lastline=15]{latex}{resources/blocks.tex}}
\only<2>{
    \centering
    \includegraphics[width=0.5\linewidth]{resources/build/blocks.pdf}\par
}
\end{frame}





\begin{frame}[fragile]
\frametitle{Diseño Multicolumna}
Para comparar datos o poner texto junto a una imagen/ecuación.

\only<1>{\inputminted[firstline=5,lastline=12]{latex}{resources/columns.tex}}
\only<2>{
    \centering
    \includegraphics[width=0.5\linewidth]{resources/build/columns.pdf}\par
}
\end{frame}





\section{Dinámica de Presentación}

\begin{frame}[fragile]
\frametitle{El comando \texttt{\textbackslash pause}}
Es la forma más simple de controlar el ritmo de la charla.

\only<1>{\inputminted[firstline=5,lastline=11]{latex}{resources/pause.tex}}
\only<2>{
    \centering
    \includegraphics[page=1,width=0.5\linewidth]{resources/build/pause.pdf}\par
}
\only<3>{
    \centering
    \includegraphics[page=2,width=0.5\linewidth]{resources/build/pause.pdf}\par
}
\only<4>{
    \centering
    \includegraphics[page=3,width=0.5\linewidth]{resources/build/pause.pdf}\par
}
\end{frame}





\begin{frame}[fragile]
\frametitle{Especificadores \texttt{< >}}
Permiten un control sofisticado sobre qué aparece y cuándo desaparece.

\only<1>{\inputminted[firstline=5,lastline=10]{latex}{resources/overlays.tex}}
\only<2>{
    \centering
    \includegraphics[page=1,width=0.5\linewidth]{resources/build/overlays.pdf}\par
}
\only<3>{
    \centering
    \includegraphics[page=2,width=0.5\linewidth]{resources/build/overlays.pdf}\par
}
\only<4>{
    \centering
    \includegraphics[page=3,width=0.5\linewidth]{resources/build/overlays.pdf}\par
}
\end{frame}





\begin{frame}
\frametitle{Enfoque Dinámico: \texttt{\textbackslash alert}}
Podemos resaltar partes del texto según la capa para dirigir la atención.

\only<1>{\inputminted[firstline=5,lastline=9]{latex}{resources/alert.tex}}
\only<2>{
    \centering
    \includegraphics[page=1,width=0.5\linewidth]{resources/build/alert.pdf}\par
}
\only<3>{
    \centering
    \includegraphics[page=2,width=0.5\linewidth]{resources/build/alert.pdf}\par
}
\only<4>{
    \centering
    \includegraphics[page=3,width=0.5\linewidth]{resources/build/alert.pdf}\par
}
\end{frame}





\section{Composición Científica}

\begin{frame}
\frametitle{Matemáticas y Fuentes}
{\lmr\beamer} usa por defecto fuentes \emph{Sans Serif}.

\only<1>{\inputminted[firstline=5,lastline=7,breaklines]{latex}{resources/equation-ss.tex}}
\only<2>{
    \centering
    \includegraphics[width=0.5\linewidth]{resources/build/equation-ss.pdf}\par
}
\end{frame}





\begin{frame}
\frametitle{Matemáticas y Fuentes}
{\lmr\beamer} para usar la fuente \emph{Serif} matemática es necesario utilizar el tema de fuentes profesionales.

\only<1>{\inputminted[firstline=3,lastline=3]{latex}{resources/equation-sf.tex}}
\only<2>{
    \centering
    \includegraphics[width=0.5\linewidth]{resources/build/equation-sf.pdf}\par
}
\end{frame}





\section{Visualización de Código}

\begin{frame}[fragile]
\frametitle{Visualización de scripts}

\only<1>{\inputminted[firstline=6,lastline=15,breaklines]{latex}{resources/minted-py.tex}}
\only<2>{
    \centering
    \includegraphics[width=0.5\linewidth]{resources/build/minted-py.pdf}\par
}
\end{frame}





\begin{frame}[fragile]
\frametitle{Visualización de scripts}

\only<1>{\inputminted[firstline=6,lastline=15,breaklines]{latex}{resources/minted-cpp.tex}}
\only<2>{
    \centering
    \includegraphics[width=0.5\linewidth]{resources/build/minted-cpp.pdf}\par
}
\end{frame}





\section{Temas y Personalización}

\begin{frame}
\frametitle{Temas de Estructura}

{\lmr\beamer} incluye temas que cambian drásticamente el look.

\pause

\begin{itemize}
    \item \textbf{Madrid:} Clásico académico, con mucha información en las barras.
    \item \textbf{Frankfurt:} Navegación superior por secciones y subsecciones.
    \item \textbf{Berkeley:} Barra lateral de navegación.
    \item \textbf{Metropolis:} Minimalista, moderno.
\end{itemize}
\end{frame}




\begin{frame}
\frametitle{Temas de Estructura}
\framesubtitle{Default}

\begin{columns}
    \begin{column}{0.5\textwidth}
        \includegraphics[page=1,width=\linewidth]{resources/build/theme-default.pdf}
    \end{column}
    \begin{column}{0.5\textwidth}
        \includegraphics[page=2,width=\linewidth]{resources/build/theme-default.pdf}
    \end{column}
\end{columns}
\end{frame}





\begin{frame}
\frametitle{Temas de Estructura}
\framesubtitle{Madrid}

\begin{columns}
    \begin{column}{0.5\textwidth}
        \includegraphics[page=1,width=\linewidth]{resources/build/theme-madrid.pdf}
    \end{column}
    \begin{column}{0.5\textwidth}
        \includegraphics[page=2,width=\linewidth]{resources/build/theme-madrid.pdf}
    \end{column}
\end{columns}
\end{frame}





\begin{frame}
\frametitle{Temas de Estructura}
\framesubtitle{Frankfurt}

\begin{columns}
    \begin{column}{0.5\textwidth}
        \includegraphics[page=1,width=\linewidth]{resources/build/theme-frankfurt.pdf}
    \end{column}
    \begin{column}{0.5\textwidth}
        \includegraphics[page=2,width=\linewidth]{resources/build/theme-frankfurt.pdf}
    \end{column}
\end{columns}
\end{frame}






\begin{frame}
\frametitle{Temas de Estructura}
\framesubtitle{Berkeley}

\begin{columns}
    \begin{column}{0.5\textwidth}
        \includegraphics[page=1,width=\linewidth]{resources/build/theme-berkeley.pdf}
    \end{column}
    \begin{column}{0.5\textwidth}
        \includegraphics[page=2,width=\linewidth]{resources/build/theme-berkeley.pdf}
    \end{column}
\end{columns}
\end{frame}





\begin{frame}
\frametitle{Temas de Estructura}
\framesubtitle{Metropolis}

\begin{columns}
    \begin{column}{0.5\textwidth}
        \includegraphics[page=1,width=\linewidth]{resources/build/theme-metropolis.pdf}
    \end{column}
    \begin{column}{0.5\textwidth}
        \includegraphics[page=3,width=\linewidth]{resources/build/theme-metropolis.pdf}
    \end{column}
\end{columns}
\end{frame}




\begin{frame}[fragile]
\frametitle{Temas de Color y Fuentes}
Podemos mezclar componentes para crear una identidad propia.
\begin{itemize}
    \item \mintinline{latex}{\usecolortheme{crane}}: Tonos amarillos y naranjas.
    \item \mintinline{latex}{\usecolortheme{beaver}}: Tonos rojos y grises.
    \item \mintinline{latex}{\usefonttheme{serif}}: Cambia todo el texto a fuentes \emph{Serif}.
\end{itemize}
\end{frame}





\begin{frame}
\frametitle{Temas de Estructura}
\framesubtitle{Madrid + Crane}

\begin{columns}
    \begin{column}{0.5\textwidth}
        \includegraphics[page=1,width=\linewidth]{resources/build/theme-madrid-crane.pdf}
    \end{column}
    \begin{column}{0.5\textwidth}
        \includegraphics[page=2,width=\linewidth]{resources/build/theme-madrid-crane.pdf}
    \end{column}
\end{columns}
\end{frame}





\begin{frame}
\frametitle{Temas de Estructura}
\framesubtitle{Madrid + Beaver}

\begin{columns}
    \begin{column}{0.5\textwidth}
        \includegraphics[page=1,width=\linewidth]{resources/build/theme-madrid-beaver.pdf}
    \end{column}
    \begin{column}{0.5\textwidth}
        \includegraphics[page=2,width=\linewidth]{resources/build/theme-madrid-beaver.pdf}
    \end{column}
\end{columns}
\end{frame}





\begin{frame}
\frametitle{Temas de Estructura}
\framesubtitle{Madrid + Serif}

\begin{columns}
    \begin{column}{0.5\textwidth}
        \includegraphics[page=1,width=\linewidth]{resources/build/theme-madrid-serif.pdf}
    \end{column}
    \begin{column}{0.5\textwidth}
        \includegraphics[page=2,width=\linewidth]{resources/build/theme-madrid-serif.pdf}
    \end{column}
\end{columns}
\end{frame}





\begin{frame}[fragile]
\frametitle{Logos}

{\lmr\beamer} también nos permite agregar un logo a la presentación de esta manera.

\begin{minted}{latex}
    \logo{\includegraphics[width=1cm]{logo.png}}
\end{minted}
\end{frame}





\section{Producción de Material}

\begin{frame}[fragile]
\frametitle{Modo \texttt{handout}}

Cuando enviamos nuestra presentación por correo, no queremos 100 páginas de \emph{overlays}.

\pause

\begin{minted}{latex}
\documentclass[handout]{beamer}
\end{minted}

\begin{itemize}
    \item Colapsa todas las capas (\emph{overlays}) en una sola página por diapositiva/\emph{frame}.
    \item Ideal para imprimir o para subir a repositorios de conferencias.
\end{itemize}
\end{frame}





\section{Buenas Prácticas}

\begin{frame}
\frametitle{Menos es Más}
\begin{itemize}
    \item No más de 6-7 líneas de texto por diapositiva.
    \item Use imágenes de alta resolución (preferiblemente vectoriales .pdf).
    \item Evite animaciones excesivas que distraigan del contenido técnico.
    \pause
    \item Una diapositiva es un apoyo al orador, no un remplazo.
\end{itemize}
\end{frame}



\section{Interactividad}

\begin{frame}[label=backup0_return]
    \frametitle{Navegación No Lineal}
    {\lmr\beamer} permite crear botones para saltar a diapositivas de respaldo.

    \vfill
    \centering
    \hyperlink{backup0}{\beamerbutton{Ver Datos de Calibración}}
    \vfill

    Útil durante la ronda de preguntas si alguien pide ver detalles técnicos omitidos.
\end{frame}





\begin{frame}[label=backup1_return]
\frametitle{Navegación No Lineal}
No necesariamente tiene que ser un botón, puede ser una gráfica o un dibujo, cualquier caja de {\lmr\LaTeX}

\vfill
\centering
\hyperlink{backup1}{
    \begin{tikzpicture}
        \begin{axis}[
            width=0.75\linewidth,
            height=0.6\paperheight,
            xlabel={Resistencia [\unit{\ohm}]},
            ylabel={Potencia [\unit{\watt}]},
        ]
            \addplot [
                red,
                mark=*,
                only marks,
                error bars,
                x dir=both,
                x explicit,
                y dir=both,
                y explicit,
            ] table [
                x error=ResistanceError,
                y error=PowerError,
            ] {resources/battery-title.csv};
        \end{axis}
    \end{tikzpicture}
}
\end{frame}





\begin{frame}[fragile]
\frametitle{Navegación No Lineal}

Este comportamiento se logra mediante el comando

\begin{minted}[escapeinside=||]{latex}
\hyperlink{|\textit{etiqueta}|}{|\textit{texto}|}
\end{minted}

Y uno de los argumentos opcional \texttt{label} del entorno \texttt{frame}, con el que asignamos esa etiqueta.
\end{frame}





\section{Inclusión de Multimedia Avanzada}

\begin{frame}
\frametitle{Multimedia}

El paquete \textbf{\texttt{multimedia}} es adecuado para incrustar audio y vídeo simple usando comandos como \mintinline{latex}{\sound} y \mintinline{latex}{\movie}.
\end{frame}



\begin{frame}
\frametitle{Multimedia}

\mintinline[escapeinside=||]{latex}{\movie[|\textit{options}|]{|\textit{poster}|}{|\textit{sample.avi}|}}
\end{frame}





\begin{frame}
\frametitle{Multimedia}

\mintinline[escapeinside=||,breaklines,breakafter=\,\[\]\{\}]{latex}{\movie[autostart,width=0.75\paperwidth,height=0.75\paperheight,loop]{}{sample.avi}}
\end{frame}





\begin{frame}
\frametitle{Multimedia}

\begin{center}
    \movie[autostart,width=0.75\paperwidth,height=0.75\paperheight,loop]{}{sample.avi}
\end{center}
\end{frame}





\begin{frame}
\frametitle{El paquete \texttt{media9}}
El paquete \texttt{media9} permite incrustar:

\begin{itemize}
    \item Vídeos de alta definición (MP4, FLV) con un reproductor SWF incrustado.
    \item Modelos 3D interactivos (U3D o PRC) esenciales en documentación científica.
\end{itemize}

Las funciones avanzadas de \texttt{media9} (3D y vídeo MP4) requieren un visor PDF con soporte para «Rich Media Annotations», como Adobe Reader.
\end{frame}





\begin{frame}[fragile]
\frametitle{Modelos 3D Interactivos}

Incrustar un modelo 3D permite a la audiencia rotar, hacer zoom y examinar la figura directamente en la diapositiva.

\pause

El archivo \texttt{.u3d} o \texttt{.prc} debe ser generado previamente con software externo (Asymptote, Meshlab, o conversión de CAD).
\end{frame}




\begin{frame}[fragile]
\frametitle{Modelos 3D Interactivos}

\begin{minted}[breaklines]{latex}
\usepackage{media9}

\includemedia[
    width=0.7\linewidth,
    height=0.5\linewidth,
    addresource=molecula.u3d,
    activate=pageopen,
    3Dcoo=1 1 1,
    3Dc2c=1,
    3Dview=default,
]{}{molecula.u3d}
\end{minted}
\end{frame}





\begin{frame}
\frametitle{Animaciones Cuadro por Cuadro con \texttt{animate}}

\pause

\begin{itemize}[<+->]
    \item El paquete \texttt{animate} se usa para crear animaciones simples a partir de una secuencia de imágenes (o código \texttt{TikZ}).
    \item El usuario puede controlar la velocidad de reproducción y pausar la animación dentro del PDF.
    \item Es perfecto para visualizar la evolución de un proceso o los pasos de un algoritmo.
\end{itemize}
\end{frame}


\begin{frame}[fragile]
\frametitle{El paquete \texttt{animate}}

\begin{minted}[breaklines]{latex}
\usepackage{animate}

\begin{animateinline}[
    autoplay,
    loop,
    poster=first,
]{10}
    \includegraphics{frame1.pdf}\newframe
    \includegraphics{frame2.pdf}\newframe
    \includegraphics{frame3.pdf}
\end{animateinline}
\end{minted}
\end{frame}





\begin{frame}[fragile]
\frametitle{El paquete \texttt{animate}}

\begin{multicols}{2}
    \scriptsize
    \begin{minted}[breaklines,breakanywhere]{latex}
\begin{animateinline}[autoplay,loop,nomouse,poster=first]{24}
    \multiframe{359}{nAngle=0+1}{
        \begin{tikzpicture}
            \draw[thick,->] (-3,0) -- (3,0) node[right] {$x$};
            \draw[thick,->] (0,-3) -- (0,3) node[above] {$y$};
            \draw[red,thick] (0,0) circle (2.5);
            \draw[ultra thick,cyan] (0,0) -- (0,0 |- \nAngle:2.5) node[midway,left] {$\sin(\alpha)$};
            \draw[ultra thick,orange] (0,0) -- (\nAngle:2.5 |- 0,0) node[midway,below] {$\cos(\alpha)$};
            \draw[densely dotted,orange] (\nAngle:2.5) -- (\nAngle:2.5 |- 0,0);
            \draw[densely dotted,cyan] (\nAngle:2.5) -- (0,0 |- \nAngle:2.5);
            \draw[thick,green] (0.5,0) arc (0:\nAngle:0.5);
            \draw[thick,green,opacity=0] (0.7,0) arc (0:\nAngle:0.7) node[midway,opacity=1] {$\alpha$};
            \draw[ultra thick,red,->] (0,0) -- (\nAngle:2.5);
        \end{tikzpicture}
    }
\end{animateinline}
    \end{minted}
\end{multicols}
\end{frame}





\begin{frame}[fragile]
\frametitle{El paquete \texttt{animate}}

\begin{center}
    \begin{animateinline}[autoplay,loop,nomouse,poster=first]{24}
        \multiframe{359}{nAngle=0+1}{
            \begin{tikzpicture}
                \draw[thick,->] (-3,0) -- (3,0) node[right] {$x$};
                \draw[thick,->] (0,-3) -- (0,3) node[above] {$y$};
                \draw[red,thick] (0,0) circle (2.5);
                \draw[ultra thick,cyan] (0,0) -- (0,0 |- \nAngle:2.5) node[midway,left] {$\sin(\alpha)$};
                \draw[ultra thick,orange] (0,0) -- (\nAngle:2.5 |- 0,0) node[midway,below] {$\cos(\alpha)$};
                \draw[densely dotted,orange] (\nAngle:2.5) -- (\nAngle:2.5 |- 0,0);
                \draw[densely dotted,cyan] (\nAngle:2.5) -- (0,0 |- \nAngle:2.5);
                \draw[thick,green] (0.5,0) arc (0:\nAngle:0.5);
                \draw[thick,green,opacity=0] (0.7,0) arc (0:\nAngle:0.7) node[midway,opacity=1] {$\alpha$};
                \draw[ultra thick,red,->] (0,0) -- (\nAngle:2.5);
            \end{tikzpicture}
        }
    \end{animateinline}
\end{center}
\end{frame}




\section{Integración Artículo-Presentación}

\begin{frame}
\frametitle{Reutilización de Contenido Científico}
\framesubtitle{Los Modos \texttt{article} y \texttt{presentation}}

\begin{itemize}[<+->]
    \item Las herramientas tradicionales (PowerPoint/Keynote) imponen una separación rígida entre el manuscrito de investigación y las diapositivas.
    \item En {\lmr\LaTeX}, existen dos clases principales de documentos:
    \begin{itemize}
        \item El modo \textbf{Artículo/Reporte} (\texttt{article}, \texttt{report}): Ideal para la redacción técnica, texto denso, bibliografía y numeración continua.
        \item El modo \textbf{Presentación} (\texttt{beamer}): Organiza el contenido en diapositivas (\texttt{frame}), con foco en visualización y aparición secuencial de elementos (\emph{overlays}).
    \end{itemize}
    \item El paquete \texttt{beamerarticle} resuelve esta duplicación al permitir que \textbf{un único archivo} se compile en cualquiera de los dos modos.
\end{itemize}
\end{frame}




\begin{frame}[fragile]
\frametitle{El Paquete \texttt{beamerarticle}}
\framesubtitle{Un código, doble propósito}

Al utilizar \mintinline{latex}{\usepackage{beamerarticle}}, {\lmr\LaTeX} introduce comandos condicionales que leen la misma fuente, pero aplican diferentes reglas de estilo.
\end{frame}





\begin{frame}[fragile]
\frametitle{El Paquete \texttt{beamerarticle}}
\framesubtitle{Un código, doble propósito}

\begin{multicols}{2}
    \scriptsize
    \inputminted[breaklines]{latex}{resources/beamerarticle-sample.tex}
\end{multicols}
\end{frame}




\begin{frame}[fragile]
\frametitle{El Paquete \texttt{beamerarticle}}
\framesubtitle{Un código, doble propósito}

\begin{center}
    \includegraphics[page=2,width=0.55\linewidth]{resources/build/beamerarticle-beamer.pdf}
\end{center}
\end{frame}





\begin{frame}[fragile]
\frametitle{El Paquete \texttt{beamerarticle}}
\framesubtitle{Un código, doble propósito}

\begin{center}
    \begin{tikzpicture}
        \node [inner sep=0,fill=white] {\includegraphics[width=0.4\linewidth]{resources/build/beamerarticle-article.pdf}};
    \end{tikzpicture}
\end{center}
\end{frame}




\begin{frame}[fragile]
\frametitle{Compilación y Modos Condicionales}
\framesubtitle{Alternando entre modos}

El control total se logra mediante la definición explícita del modo de compilación en el comando \mintinline{latex}{\documentclass}.

\bigskip

\begin{columns}
    \begin{column}{0.5\textwidth}
        \textbf{\Large Modo Artículo}
        \bigskip
        \begin{minted}{latex}
\documentclass{article}
\usepackage{beamerarticle}
        \end{minted}
    \end{column}
    \begin{column}{0.5\textwidth}
        \textbf{\Large Modo Presentación}
        \bigskip
        \begin{minted}{latex}
        \documentclass{beamer}
        \end{minted}
    \end{column}
\end{columns}
\end{frame}





\begin{frame}
\frametitle{Compilación y Modos Condicionales}
\framesubtitle{Modos Condicionales}


El comando \mintinline[escapeinside=||]{latex}{\mode<|\textit{modo}|>{|\textit{código}|}} para que un bloque se ejecute solo en el salida deseada.

Este modificador funciona con todos los comandos a los que se les puede aplicar un \emph{overlay}.
\end{frame}





\begin{frame}[allowframebreaks]
\frametitle{Referencias}
\nocite{*}
\printbibliography
\end{frame}


\begin{frame}[label=backup0]
    \frametitle{Datos de Calibración}
    Esta diapositiva está oculta en el flujo normal, pero es accesible vía el botón anterior.
    \begin{itemize}
        \item Detector: Criogénico HPGe.
        \item Tiempo de integración: 48 horas.
        \item Resolución: 1.8 keV a 1.33 MeV.
    \end{itemize}
    \vfill\hfill
    \hyperlink{backup0_return}{\beamerreturnbutton{Volver a la presentación}}
\end{frame}


\begin{frame}[label=backup1]
    \frametitle{Datos de la Gráfica}
    \begin{center}
        \tiny
        \begin{tabular}{S[table-format=2.1(1)] S[table-format=0.6(1)]}
            \toprule
            {Resistencia [\unit{\ohm}]} & {Potencia [\unit{\watt}]} \\
            \midrule
            50.0 +- 2.5 & 0.036990 +- 0.007039 \\
            47.0 +- 2.5 & 0.036990 +- 0.007039 \\
            44.0 +- 2.5 & 0.038360 +- 0.007053 \\
            41.0 +- 2.5 & 0.042160 +- 0.007050 \\
            38.0 +- 2.5 & 0.043200 +- 0.007018 \\
            35.0 +- 2.5 & 0.046900 +- 0.007021 \\
            32.0 +- 2.5 & 0.048840 +- 0.006963 \\
            29.0 +- 2.5 & 0.052000 +- 0.006929 \\
            26.0 +- 2.5 & 0.055470 +- 0.006947 \\
            23.0 +- 2.5 & 0.060000 +- 0.006882 \\
            20.0 +- 2.5 & 0.064660 +- 0.006879 \\
            17.0 +- 2.5 & 0.069600 +- 0.006936 \\
            14.0 +- 2.5 & 0.079800 +- 0.007080 \\
            11.0 +- 2.5 & 0.086920 +- 0.007232 \\
            8.0 +- 2.5 & 0.085500 +- 0.007262 \\
            5.0 +- 2.5 & 0.101260 +- 0.008415 \\
            2.0 +- 2.5 & 0.091000 +- 0.011203 \\
            0.0 +- 2.5 & 0.081480 +- 0.011828 \\
            \bottomrule
        \end{tabular}
    \end{center}
    \vfill\hfill
    \hyperlink{backup1_return}{\beamerreturnbutton{Volver a la presentación}}
\end{frame}
\end{document}
